\documentclass[11pt,a4paper]{article}

\usepackage[utf8]{inputenc}
\usepackage[T1]{fontenc}
\usepackage[margin=2.5cm]{geometry}
\usepackage{setspace}
\usepackage{lmodern}
\usepackage{microtype}
\usepackage{xcolor}
\usepackage{titlesec}
\usepackage{hyperref}
\usepackage{bookmark}
\usepackage{enumitem}
\usepackage{amsmath}
\usepackage{listings}
\usepackage[most]{tcolorbox}
\usepackage{fancyhdr}

\tcbuselibrary{listings,skins,breakable}

% -------------------------
% Corporate palette
% -------------------------
\definecolor{PrimaryNavy}{HTML}{0B1F3A}
\definecolor{AccentBlue}{HTML}{1E5AA6}
\definecolor{BodyGray}{HTML}{2F3640}
\definecolor{MutedGray}{HTML}{6B7280}
\definecolor{SoftBlue}{HTML}{EEF4FF}
\definecolor{CodeBg}{HTML}{F6F8FA}
\definecolor{CodeFrame}{HTML}{D9DEE8}
\definecolor{CodeKeyword}{HTML}{0B5394}
\definecolor{CodeString}{HTML}{38761D}
\definecolor{CodeComment}{HTML}{7A7A7A}
\definecolor{CodeNumber}{HTML}{AA3731}

% -------------------------
% Global style
% -------------------------
\setstretch{1.18}
\setlength{\parindent}{0pt}
\setlength{\parskip}{0.55em}

\AtBeginDocument{\color{BodyGray}}

\titleformat{\section}
  {\Large\bfseries\color{PrimaryNavy}}
  {\thesection}{0.7em}{}
\titleformat{\subsection}
  {\large\bfseries\color{AccentBlue}}
  {\thesubsection}{0.7em}{}
\titleformat{\subsubsection}
  {\normalsize\bfseries\color{PrimaryNavy}}
  {\thesubsubsection}{0.7em}{}

\titlespacing*{\section}{0pt}{2.0ex plus 0.3ex}{1.1ex}
\titlespacing*{\subsection}{0pt}{1.6ex plus 0.3ex}{0.7ex}

\hypersetup{
  colorlinks=true,
  linkcolor=AccentBlue,
  urlcolor=AccentBlue,
  citecolor=AccentBlue,
  pdfborder={0 0 0},
  pdftitle={Enterprise Architecture Blueprint},
  pdfauthor={Architecture Whitepaper}
}

\pagestyle{fancy}
\fancyhf{}
\fancyhead[L]{\textcolor{MutedGray}{Enterprise Architecture Blueprint}}
\fancyhead[R]{\textcolor{MutedGray}{\thepage}}
\renewcommand{\headrulewidth}{0.2pt}

% -------------------------
% Listings and boxes
% -------------------------
\lstdefinestyle{enterprise}{
  backgroundcolor=\color{CodeBg},
  basicstyle=\ttfamily\small\color{BodyGray},
  keywordstyle=\color{CodeKeyword}\bfseries,
  stringstyle=\color{CodeString},
  commentstyle=\color{CodeComment}\itshape,
  numberstyle=\scriptsize\color{MutedGray},
  numbers=left,
  numbersep=8pt,
  breaklines=true,
  showstringspaces=false,
  tabsize=2,
  frame=none,
  keepspaces=true
}

\newtcblisting{codeblock}[2][]{
  enhanced,
  breakable,
  colback=CodeBg,
  colframe=CodeFrame,
  boxrule=0.6pt,
  arc=2.8mm,
  left=1mm,
  right=1mm,
  top=1mm,
  bottom=1mm,
  title={#2},
  fonttitle=\bfseries\color{PrimaryNavy},
  listing only,
  listing options={style=enterprise,#1}
}

\newtcolorbox{conceptbox}{
  enhanced,
  breakable,
  colback=SoftBlue,
  colframe=AccentBlue!55,
  boxrule=0.7pt,
  arc=3mm,
  left=2.2mm,
  right=2.2mm,
  top=1.5mm,
  bottom=1.5mm
}

\begin{document}

\begin{titlepage}
  \centering
  \vspace*{2.8cm}

  {\Huge\bfseries\color{PrimaryNavy} Enterprise Architecture Blueprint\par}
  \vspace{0.55cm}
  {\Large\color{AccentBlue} Beginner-Friendly + Production-Grade\par}
  \vspace{0.2cm}
  {\Large\color{AccentBlue} Preventive Audit Platform Whitepaper\par}

  \vspace{1.2cm}
  {\color{MutedGray}\rule{0.78\textwidth}{0.6pt}\par}

  \vfill
  {\large\color{BodyGray}Whitepaper Edition\par}
  \vspace{0.35cm}
  {\large\color{BodyGray}\today\par}

  \vspace*{1.5cm}
\end{titlepage}

\tableofcontents
\newpage

\section{Phase 1: The Bird's-Eye View (System Overview for Dummies)}

\subsection{Conceptual explanation (for beginners)}
\begin{conceptbox}
Public procurement data is large, noisy, and constantly changing. If analysis only uses a current snapshot, high-risk behavioral signals are often missed, for example:
\begin{itemize}[leftmargin=1.4em]
  \item Founders changing immediately before a tender;
  \item ``Fake competition'' among related suppliers;
  \item Abnormal rejection patterns or suspicious pricing dynamics;
  \item Technical specifications that appear crafted for one vendor.
\end{itemize}

The platform is designed as a preventive audit system that:
\begin{itemize}[leftmargin=1.4em]
  \item Collects source data (API first, scraping fallback),
  \item Stores immutable raw evidence,
  \item Cleans and links entities (companies, people, addresses),
  \item Runs explicit and explainable risk rules (no black-box ML),
  \item Generates legal-safe outputs (PDF dossier + DOCX complaint draft).
\end{itemize}
\end{conceptbox}

\subsection{Technical design}
\textbf{High-level data flow:}
\begin{quote}
eGov/Goszakup/KGD sources $\rightarrow$ Ingestion $\rightarrow$ Raw storage $\rightarrow$ Validation $\rightarrow$ Core canonical data $\rightarrow$ Outbox events $\rightarrow$ Risk engine $\rightarrow$ Alerts $\rightarrow$ Legal dossier + complaint draft
\end{quote}

\textbf{Layered architecture:}
\begin{itemize}[leftmargin=1.4em]
  \item \textbf{Collection plane:} API clients + browser automation for dynamic pages.
  \item \textbf{Ingestion plane:} immutable write path for raw JSON/PDF metadata.
  \item \textbf{Normalization plane:} validation and idempotent upserts of canonical entities.
  \item \textbf{Event plane:} outbox + change queue.
  \item \textbf{Analytics plane:} rule engine + graph analysis + geospatial analysis.
  \item \textbf{Presentation plane:} API and legal document generation.
\end{itemize}

\textbf{Core principles:}
\begin{itemize}[leftmargin=1.4em]
  \item API-first, scraping as fallback,
  \item Raw data as immutable evidence,
  \item No business logic in database triggers,
  \item Idempotent pipeline with safe retries,
  \item OLTP and OLAP separation,
  \item Explainable scoring for legal defensibility.
\end{itemize}

\subsection{Key snippets}
\begin{codeblock}[language={}]{System flow (text diagram)}
[Data Sources]
  |- Goszakup GraphQL/API
  |- eGov open data
  |- KGD files (CSV/XML/PDF)
  v
[Ingest Workers]
  -> raw.goszakup_protocols (JSONB)
  -> raw.documents_registry
  v
[Application Service]
  -> validate (Pydantic v2)
  -> core upsert (company_dim, procurement_protocol_fact)
  -> ops.dead_letter (bad records)
  -> ops.outbox_events (domain events)
  v
[Outbox Relay]
  -> ops.change_queue
  v
[Scoring Dispatcher]
  -> rule strategies
  -> scoring.alerts + scoring.alert_explanations
  v
[FastAPI + Document Factory]
  -> Alert UI
  -> PDF dossier
  -> DOCX complaint draft for eOtinish
\end{codeblock}

\section{Phase 2: Database Architecture (CQRS Pattern --- OLTP vs OLAP)}

\subsection{Conceptual explanation (for beginners)}
\begin{conceptbox}
A single database node cannot efficiently serve continuous ingestion writes and heavy analytics at the same time. OLTP and OLAP workloads compete for I/O, CPU, locks, and memory.

When analytical joins and recomputations run on a write-heavy node, common failure modes appear:
\begin{itemize}[leftmargin=1.4em]
  \item I/O spikes and unstable write latency,
  \item lock contention,
  \item checkpoint pressure,
  \item delayed ingestion and unreliable SLAs.
\end{itemize}

CQRS in this context means separating write truth from read/analytics processing:
\begin{itemize}[leftmargin=1.4em]
  \item \textbf{Master node (OLTP):} write-fast, transactional integrity, minimal heavy reads.
  \item \textbf{Analytics node (OLAP):} heavy joins, geospatial logic, temporal rollups.
\end{itemize}
\end{conceptbox}

\subsection{Technical design}
\textbf{Master node (PostgreSQL 15, OLTP):}
\begin{itemize}[leftmargin=1.4em]
  \item Purpose: high-throughput ingestion of raw JSON and canonical transactional persistence.
  \item Schemas: \texttt{raw}, \texttt{core}, \texttt{ops}.
  \item Partitioning by date where beneficial.
  \item Minimal expensive indexes on raw ingestion tables.
  \item Strict idempotency keys (\texttt{source\_cursor}, \texttt{checksum}, natural business keys).
\end{itemize}

\textbf{Analytics node (OLAP):}
\begin{itemize}[leftmargin=1.4em]
  \item Replication from master (logical replication/CDC preferred).
  \item Extensions: \texttt{postgis} for spatial rules, \texttt{timescaledb} for temporal analytics.
  \item Additional schemas: \texttt{mart}, \texttt{scoring}, graph export staging.
  \item Nightly refresh windows for heavy rollups/materialized views.
\end{itemize}

\subsection{Key snippets}
\begin{codeblock}[language=SQL]{SQL bootstrapping}
CREATE SCHEMA IF NOT EXISTS raw;
CREATE SCHEMA IF NOT EXISTS core;
CREATE SCHEMA IF NOT EXISTS ops;
CREATE SCHEMA IF NOT EXISTS mart;
CREATE SCHEMA IF NOT EXISTS scoring;
\end{codeblock}

\begin{codeblock}[language=SQL]{OLTP critical tables}
CREATE TABLE IF NOT EXISTS raw.goszakup_protocols (
    id BIGSERIAL PRIMARY KEY,
    source_cursor TEXT NOT NULL,
    source_batch_id UUID NOT NULL,
    checksum TEXT NOT NULL,
    payload JSONB NOT NULL,
    scraped_at TIMESTAMPTZ NOT NULL DEFAULT now(),
    UNIQUE (source_cursor, checksum)
);

CREATE TABLE IF NOT EXISTS core.company_dim (
    company_id BIGSERIAL PRIMARY KEY,
    company_bin TEXT NOT NULL UNIQUE,
    company_name TEXT,
    company_type TEXT,
    address_raw TEXT,
    address_norm TEXT,
    rka_code TEXT,
    geog GEOGRAPHY(POINT, 4326),
    source_raw_id BIGINT,
    source_loaded_at TIMESTAMPTZ NOT NULL DEFAULT now(),
    updated_at TIMESTAMPTZ NOT NULL DEFAULT now()
);

CREATE TABLE IF NOT EXISTS core.procurement_protocol_fact (
    protocol_fact_id BIGSERIAL PRIMARY KEY,
    protocol_id TEXT NOT NULL,
    lot_id TEXT NOT NULL,
    protocol_date DATE NOT NULL,
    buyer_bin TEXT,
    supplier_bin TEXT,
    amount NUMERIC(18,2),
    protocol_status TEXT,
    rejection_reason TEXT,
    is_antidumping BOOLEAN,
    source_raw_id BIGINT,
    source_loaded_at TIMESTAMPTZ NOT NULL DEFAULT now(),
    updated_at TIMESTAMPTZ NOT NULL DEFAULT now(),
    UNIQUE (protocol_date, protocol_id, lot_id)
);
\end{codeblock}

\begin{codeblock}[language={}]{postgresql.conf (write-heavy OLTP baseline)}
# Memory / connections
max_connections = 200
shared_buffers = 8GB
effective_cache_size = 24GB
work_mem = 16MB
maintenance_work_mem = 1GB

# WAL / checkpoints (write-heavy)
wal_level = logical
synchronous_commit = local
wal_compression = on
max_wal_size = 16GB
min_wal_size = 4GB
checkpoint_timeout = 30min
checkpoint_completion_target = 0.9

# Parallelism / I/O
max_worker_processes = 16
max_parallel_workers = 8
max_parallel_workers_per_gather = 2
effective_io_concurrency = 200
random_page_cost = 1.1

# Autovacuum tuned for JSON-heavy writes
autovacuum = on
autovacuum_naptime = 10s
autovacuum_vacuum_scale_factor = 0.02
autovacuum_analyze_scale_factor = 0.01
\end{codeblock}

\begin{codeblock}[language=SQL]{OLAP extension setup}
CREATE EXTENSION IF NOT EXISTS postgis;
CREATE EXTENSION IF NOT EXISTS timescaledb;

CREATE TABLE IF NOT EXISTS mart.buyer_daily_metrics (
    day DATE NOT NULL,
    buyer_bin TEXT NOT NULL,
    total_tenders INT NOT NULL,
    rejected_count INT NOT NULL,
    PRIMARY KEY (day, buyer_bin)
);

SELECT create_hypertable(
    'mart.buyer_daily_metrics',
    by_range('day'),
    if_not_exists => TRUE
);
\end{codeblock}

\section{Phase 3: Data Ingestion \& Evasion Layer (Scraping \& OCR)}

\subsection{Conceptual explanation (for beginners)}
\begin{conceptbox}
Public-sector sources are heterogeneous:
\begin{itemize}[leftmargin=1.4em]
  \item some records are available through APIs,
  \item some are only visible in dynamic JavaScript pages,
  \item some are uploaded as scanned PDFs.
\end{itemize}

Therefore ingestion must be resilient and staged, not a single fragile script:
\begin{enumerate}[leftmargin=1.5em]
  \item Fetch legally available source,
  \item Preserve exact raw payload,
  \item Extract linked documents,
  \item OCR image-only documents,
  \item Normalize noisy text,
  \item Send cleaned record to validation and upsert.
\end{enumerate}
\end{conceptbox}

\subsection{Technical design}
\textbf{Access strategy (compliance-first):}
\begin{itemize}[leftmargin=1.4em]
  \item Priority 1: official APIs and published datasets.
  \item Priority 2: browser automation for data unavailable via API.
  \item Reliability controls: rate limits, retries, request budgets, rotating egress pools.
  \item Source audit log per request for traceability.
\end{itemize}

\textbf{Headless browser tier:}
\begin{itemize}[leftmargin=1.4em]
  \item Playwright workers with user-agent rotation.
  \item Session/cookie isolation.
  \item Retry backoff and circuit breaker.
  \item Proxy pool abstraction (residential providers if contractually/legal approved).
  \item CAPTCHA path via manual queue or approved service integration.
\end{itemize}

\textbf{OCR pipeline:}
\begin{itemize}[leftmargin=1.4em]
  \item Detect text-PDF vs scanned PDF.
  \item If text-PDF: direct extraction.
  \item If scanned: page render + OCR with \texttt{kaz+rus+eng}.
  \item Persist OCR text, confidence score, and page-level metadata.
  \item Route low-confidence pages to manual review.
\end{itemize}

\textbf{Dirty data normalization:}
\begin{itemize}[leftmargin=1.4em]
  \item Unicode normalization (\texttt{NFKC}).
  \item Cyrillic/Latin confusable replacement.
  \item Address standardization + RKA geocoding where possible.
  \item Number/currency/date normalization.
\end{itemize}

\subsection{Key snippets}
\begin{codeblock}[language=Python]{Playwright worker skeleton (Python)}
from playwright.async_api import async_playwright

async def fetch_dynamic_page(url: str, proxy: dict | None):
    async with async_playwright() as p:
        browser = await p.chromium.launch(headless=True, proxy=proxy)
        ctx = await browser.new_context(user_agent="Mozilla/5.0 ...")
        page = await ctx.new_page()
        await page.goto(url, wait_until="networkidle", timeout=60_000)
        content = await page.content()
        await browser.close()
        return content
\end{codeblock}

\begin{codeblock}[language=Python]{OCR detection and extraction}
import fitz  # PyMuPDF
import pytesseract
from PIL import Image

def pdf_has_text(pdf_path: str) -> bool:
    doc = fitz.open(pdf_path)
    for page in doc:
        if page.get_text().strip():
            return True
    return False

def ocr_pdf_page_images(images: list[Image.Image]) -> str:
    chunks = []
    for i, img in enumerate(images):
        text = pytesseract.image_to_string(img, lang="kaz+rus+eng")
        chunks.append(f"\n--- page {i+1} ---\n{text}")
    return "".join(chunks)
\end{codeblock}

\begin{codeblock}[language=Python]{Confusable Cyrillic/Latin cleanup}
CONFUSABLE_MAP = str.maketrans({
    "A": "А", "B": "В", "C": "С", "E": "Е", "H": "Н", "K": "К",
    "M": "М", "O": "О", "P": "Р", "T": "Т", "X": "Х",
    "a": "а", "c": "с", "e": "е", "o": "о", "p": "р", "x": "х", "y": "у"
})

def normalize_text(s: str) -> str:
    import unicodedata
    import re
    s = unicodedata.normalize("NFKC", s)
    s = s.translate(CONFUSABLE_MAP)
    s = re.sub(r"\s+", " ", s).strip()
    return s
\end{codeblock}

\section{Phase 4: Application Layer \& Transactional Pipeline (Outbox Pattern)}

\subsection{Conceptual explanation (for beginners)}
\begin{conceptbox}
The dual-write problem occurs when an application:
\begin{enumerate}[leftmargin=1.5em]
  \item writes business data to a database, then
  \item publishes an event to a queue,
\end{enumerate}
and crashes between these actions.

Outbox pattern addresses this by writing business updates and event rows in the same database transaction. A separate relay worker later publishes pending outbox events to the queue.

This guarantees two critical properties:
\begin{itemize}[leftmargin=1.4em]
  \item no lost events,
  \item no phantom events.
\end{itemize}
\end{conceptbox}

\subsection{Technical design}
\textbf{Pipeline components:}
\begin{itemize}[leftmargin=1.4em]
  \item \texttt{RawBatchReader} (watermark + batch leasing),
  \item \texttt{RawRecordValidator} (Pydantic v2),
  \item \texttt{CoreUpsertService} (idempotent upserts, no triggers),
  \item \texttt{OutboxWriter} (domain events),
  \item \texttt{OutboxRelayWorker} (\texttt{ops.outbox\_events} to \texttt{ops.change\_queue}),
  \item \texttt{DeadLetterWriter} (\texttt{ops.dead\_letter}),
  \item \texttt{UnitOfWork} with explicit transaction boundaries.
\end{itemize}

\textbf{Idempotency keys:}
\begin{itemize}[leftmargin=1.4em]
  \item Raw layer: (\texttt{source\_cursor}, \texttt{checksum}),
  \item Fact layer: (\texttt{protocol\_date}, \texttt{protocol\_id}, \texttt{lot\_id}),
  \item Company layer: \texttt{company\_bin},
  \item Alert/event layer: fingerprint + UUID.
\end{itemize}

\subsection{Key snippets}
\begin{codeblock}[language=SQL]{Outbox + dead-letter DDL}
CREATE TABLE IF NOT EXISTS ops.dead_letter (
    dead_letter_id BIGSERIAL PRIMARY KEY,
    raw_id BIGINT,
    source_cursor TEXT,
    source_batch_id UUID,
    checksum TEXT,
    error_type TEXT NOT NULL,
    error_message TEXT NOT NULL,
    validation_errors JSONB,
    payload JSONB,
    created_at TIMESTAMPTZ NOT NULL DEFAULT now()
);

CREATE TABLE IF NOT EXISTS ops.outbox_events (
    event_id UUID PRIMARY KEY,
    event_type TEXT NOT NULL,
    aggregate_type TEXT NOT NULL,
    aggregate_id TEXT NOT NULL,
    payload JSONB NOT NULL,
    occurred_at TIMESTAMPTZ NOT NULL DEFAULT now(),
    published_at TIMESTAMPTZ,
    publish_status TEXT NOT NULL DEFAULT 'pending',
    source_batch_id UUID
);

CREATE TABLE IF NOT EXISTS ops.change_queue (
    queue_id BIGSERIAL PRIMARY KEY,
    entity_type TEXT NOT NULL,
    entity_id TEXT NOT NULL,
    bucket_date DATE,
    payload JSONB,
    status TEXT NOT NULL DEFAULT 'pending',
    attempts INT NOT NULL DEFAULT 0,
    created_at TIMESTAMPTZ NOT NULL DEFAULT now(),
    updated_at TIMESTAMPTZ NOT NULL DEFAULT now()
);
\end{codeblock}

\begin{codeblock}[language=SQL]{Idempotent upsert (SQL)}
INSERT INTO core.procurement_protocol_fact (
    protocol_id, lot_id, protocol_date, buyer_bin, supplier_bin,
    amount, protocol_status, rejection_reason, is_antidumping,
    source_raw_id, source_loaded_at
)
VALUES (
    :protocol_id, :lot_id, :protocol_date, :buyer_bin, :supplier_bin,
    :amount, :protocol_status, :rejection_reason, :is_antidumping,
    :source_raw_id, now()
)
ON CONFLICT (protocol_date, protocol_id, lot_id)
DO UPDATE SET
    buyer_bin = EXCLUDED.buyer_bin,
    supplier_bin = EXCLUDED.supplier_bin,
    amount = EXCLUDED.amount,
    protocol_status = EXCLUDED.protocol_status,
    rejection_reason = EXCLUDED.rejection_reason,
    is_antidumping = EXCLUDED.is_antidumping,
    source_raw_id = EXCLUDED.source_raw_id,
    source_loaded_at = now(),
    updated_at = now();
\end{codeblock}

\begin{codeblock}[language=Python]{SQLAlchemy 2.0 transactional unit}
from sqlalchemy.orm import Session
from uuid import uuid4

def process_raw_record(session: Session, raw_rec: dict):
    try:
        # 1) Validate payload
        model = RawGoszakupProtocolPayload.model_validate(raw_rec["payload"])

        with session.begin():
            # 2) Upsert canonical entities
            upsert_company(session, model.buyer)
            upsert_company(session, model.supplier)
            upsert_protocol_fact(session, model, source_raw_id=raw_rec["id"])

            # 3) Write outbox event in the same transaction
            event = {
                "event_id": str(uuid4()),
                "event_type": "ProcurementProtocolUpserted",
                "aggregate_type": "ProcurementProtocol",
                "aggregate_id": f"{model.protocol_id}:{model.lot_id}",
                "payload": model.model_dump(mode="json"),
                "source_batch_id": raw_rec["source_batch_id"]
            }
            insert_outbox(session, event)
    except Exception as exc:
        write_dead_letter(session, raw_rec, exc)
\end{codeblock}

\begin{codeblock}[language=SQL]{Relay query pattern}
SELECT event_id, event_type, aggregate_type, aggregate_id, payload
FROM ops.outbox_events
WHERE publish_status = 'pending'
ORDER BY occurred_at
FOR UPDATE SKIP LOCKED
LIMIT 500;
\end{codeblock}

\section{Phase 5: Historical Tracking \& Graph Database (SCD Type 2 \& Neo4j)}

\subsection{Conceptual explanation (for beginners)}
\begin{conceptbox}
If company founders or addresses are overwritten in place, historical evidence is destroyed. Many procurement risks are temporal by nature (for example, a founder change shortly before tender publication).

SCD Type 2 preserves full history by storing versions with:
\begin{itemize}[leftmargin=1.4em]
  \item \texttt{valid\_from},
  \item \texttt{valid\_to} (NULL for active version),
  \item \texttt{is\_current}.
\end{itemize}

This enables time-accurate questions such as:
\begin{quote}
``Who was the founder on the tender date?''
\end{quote}
\end{conceptbox}

\subsection{Technical design}
\textbf{Relational temporal layer (SCD2):}
\begin{itemize}[leftmargin=1.4em]
  \item \texttt{core.company\_dim\_scd}
  \item \texttt{core.person\_dim\_scd}
  \item \texttt{core.company\_person\_rel\_scd}
\end{itemize}

\textbf{Update semantics:}
\begin{enumerate}[leftmargin=1.5em]
  \item close current row (\texttt{valid\_to = now()}, \texttt{is\_current = false}),
  \item insert a new current row with changed attributes.
\end{enumerate}

\textbf{Neo4j temporal graph model:}
\begin{itemize}[leftmargin=1.4em]
  \item Nodes: \texttt{Company}, \texttt{Person}, \texttt{Address}, \texttt{TenderLot}.
  \item Time-valid edges:
  \begin{itemize}[leftmargin=1.4em]
    \item \texttt{(:Person)-[:IS\_FOUNDER\_OF \{valid\_from, valid\_to\}]->(:Company)}
    \item \texttt{(:Person)-[:IS\_DIRECTOR\_OF \{...\}]->(:Company)}
    \item \texttt{(:Company)-[:REGISTERED\_AT \{...\}]->(:Address)}
    \item \texttt{(:Company)-[:PARTICIPATED\_IN \{lot\_id, date\}]->(:TenderLot)}
  \end{itemize}
\end{itemize}

\textbf{Entity resolution:}
\begin{itemize}[leftmargin=1.4em]
  \item Deterministic matching by BIN/IIN first,
  \item Probabilistic matching for name/address variants,
  \item Phonetic normalization and RKA-aware address geocoding,
  \item Persist \texttt{confidence\_score} and \texttt{match\_method}.
\end{itemize}

\subsection{Key snippets}
\begin{codeblock}[language=SQL]{SCD Type 2 DDL}
CREATE TABLE IF NOT EXISTS core.company_dim_scd (
    company_scd_id BIGSERIAL PRIMARY KEY,
    company_bin TEXT NOT NULL,
    company_name TEXT,
    address_norm TEXT,
    rka_code TEXT,
    company_type TEXT,
    valid_from TIMESTAMPTZ NOT NULL,
    valid_to TIMESTAMPTZ,
    is_current BOOLEAN NOT NULL DEFAULT TRUE,
    hash_diff TEXT NOT NULL
);

CREATE UNIQUE INDEX IF NOT EXISTS ux_company_current
ON core.company_dim_scd(company_bin)
WHERE is_current = TRUE;
\end{codeblock}

\begin{codeblock}[language=SQL]{SCD2 update transaction (concept)}
-- Close current
UPDATE core.company_dim_scd
SET valid_to = now(), is_current = FALSE
WHERE company_bin = :bin
  AND is_current = TRUE
  AND hash_diff <> :new_hash;

-- Open new version if changed
INSERT INTO core.company_dim_scd (
  company_bin, company_name, address_norm, rka_code, company_type,
  valid_from, valid_to, is_current, hash_diff
)
SELECT :bin, :name, :addr, :rka, :ctype, now(), NULL, TRUE, :new_hash
WHERE EXISTS (
  SELECT 1
  WHERE NOT EXISTS (
    SELECT 1
    FROM core.company_dim_scd
    WHERE company_bin = :bin
      AND is_current = TRUE
      AND hash_diff = :new_hash
  )
);
\end{codeblock}

\begin{codeblock}[language={}]{Neo4j constraints + temporal edge merge}
CREATE CONSTRAINT company_bin IF NOT EXISTS
FOR (c:Company) REQUIRE c.bin IS UNIQUE;

CREATE CONSTRAINT person_iin IF NOT EXISTS
FOR (p:Person) REQUIRE p.iin IS UNIQUE;

MERGE (p:Person {iin: $iin})
MERGE (c:Company {bin: $bin})
MERGE (p)-[r:IS_FOUNDER_OF {source_ref: $source_ref}]->(c)
SET r.valid_from = datetime($valid_from),
    r.valid_to   = datetime($valid_to);
\end{codeblock}

\section{Phase 6: The Risk Scoring Engine (Rule-Based, No ML)}

\subsection{Conceptual explanation (for beginners)}
\begin{conceptbox}
At this stage, deterministic rule-based scoring is preferred to machine learning because legal defensibility and traceability are mandatory.

Reasons to avoid black-box models in initial production:
\begin{itemize}[leftmargin=1.4em]
  \item difficult to explain in legal or regulatory review,
  \item potentially high false-positive cost,
  \item weak line-by-line evidence mapping.
\end{itemize}

Rule-based scoring provides:
\begin{itemize}[leftmargin=1.4em]
  \item explicit indicators,
  \item transparent scores,
  \item explainable evidence for each alert.
\end{itemize}
\end{conceptbox}

\subsection{Technical design}
\textbf{Strategy Pattern components:}
\begin{itemize}[leftmargin=1.4em]
  \item \texttt{BaseScoringRule},
  \item concrete rules:
    \texttt{AkimatKillerRule},
    \texttt{AntiDumpingSniperRule},
    \texttt{SpatialAnomalyRule},
    \texttt{HiddenNetworkRule}.
\end{itemize}

\textbf{Dispatcher flow:}
\begin{enumerate}[leftmargin=1.5em]
  \item lock pending item from \texttt{ops.change\_queue} with \texttt{FOR UPDATE SKIP LOCKED},
  \item select applicable rules by \texttt{entity\_type},
  \item execute rules and compute candidates,
  \item upsert alerts and explanations,
  \item mark queue item done/failed and write \texttt{ops.rule\_run\_log}.
\end{enumerate}

\textbf{Risk formula:}
\[
RI(x) = 100 \cdot \sum_{i=1}^{n} w_i \cdot I_i(x), \qquad
\sum_{i=1}^{n} w_i = 1, \qquad
I_i(x) \in [0,1].
\]

Store weights and thresholds in a DB/UI configuration table so calibration does not require redeployment.

\subsection{Key snippets}
\begin{codeblock}[language=Python]{Strategy interface}
from dataclasses import dataclass
from typing import Protocol, List, Dict, Any

@dataclass
class AlertCandidate:
    rule_name: str
    rule_version: int
    entity_type: str
    entity_id: str
    score: float
    severity: str
    status: str
    window_start: str
    window_end: str
    fingerprint: str
    evidence: Dict[str, Any]

class BaseScoringRule(Protocol):
    rule_name: str
    rule_version: int
    def applicable_entities(self) -> list[str]: ...
    def compute(self, context) -> List[AlertCandidate]: ...
\end{codeblock}

\begin{codeblock}[language=SQL]{Hidden network idea (SQL CTE with shared identifiers)}
WITH lot_parties AS (
  SELECT lot_id, supplier_bin
  FROM core.procurement_protocol_fact
  WHERE protocol_date >= current_date - INTERVAL '180 days'
),
ident AS (
  SELECT company_bin, lower(email) AS email, regexp_replace(phone, '\D','','g') AS phone
  FROM core.company_contacts
),
shared AS (
  SELECT a.lot_id,
         i1.company_bin AS c1,
         i2.company_bin AS c2,
         (i1.email = i2.email AND i1.email IS NOT NULL) AS shared_email,
         (i1.phone = i2.phone AND i1.phone IS NOT NULL) AS shared_phone
  FROM lot_parties a
  JOIN lot_parties b
    ON a.lot_id = b.lot_id
   AND a.supplier_bin < b.supplier_bin
  JOIN ident i1 ON i1.company_bin = a.supplier_bin
  JOIN ident i2 ON i2.company_bin = b.supplier_bin
)
SELECT *
FROM shared
WHERE shared_email OR shared_phone;
\end{codeblock}

\begin{codeblock}[language={}]{Betweenness centrality example (Neo4j GDS)}
CALL gds.betweenness.stream('procurementGraph')
YIELD nodeId, score
RETURN gds.util.asNode(nodeId).bin AS company_bin, score
ORDER BY score DESC
LIMIT 50;
\end{codeblock}

\section{Phase 7: Legal Engineering \& Document Generation (The Output)}

\subsection{Conceptual explanation (for beginners)}
\begin{conceptbox}
The system must produce legally usable outputs, not only anomaly flags.

Primary end-user flow (auditor/legal expert):
\begin{enumerate}[leftmargin=1.5em]
  \item Open alert list,
  \item Filter by severity/rule/region,
  \item Open a case and inspect timeline/evidence,
  \item Generate both:
    \begin{itemize}[leftmargin=1.4em]
      \item visual PDF dossier,
      \item DOCX complaint draft ready for eOtinish submission.
    \end{itemize}
\end{enumerate}
\end{conceptbox}

\subsection{Technical design}
\textbf{UX/UI flow:}
\begin{itemize}[leftmargin=1.4em]
  \item Alerts list with traffic-light severity coding.
  \item Alert detail with indicator-level score breakdown.
  \item Evidence blocks + linked entity graph + chronology.
  \item ``Generate Documents'' action:
    completeness validation,
    linguistic shield transform,
    PDF/DOCX render,
    artifact audit logging.
\end{itemize}

\textbf{Linguistic shield policy:}
\begin{itemize}[leftmargin=1.4em]
  \item Replace legally aggressive statements with neutral, probabilistic formulations.
  \item Always append disclaimer:
  \begin{quote}
  ``Algorithmic risk assessment based on open data; not a direct assertion of criminal guilt.''
  \end{quote}
\end{itemize}

\subsection{Key snippets}
\begin{codeblock}[language=Python]{Linguistic shield dictionary}
SHIELD_TERMS = {
    "картельный сговор": "признаки синхронности поведения и потенциальной аффилированности",
    "фиктивное ТОО": "компания с признаками высокой транзитивности",
    "заточенный тендер": "техническая спецификация с аномально узкими критериями"
}

def apply_linguistic_shield(text: str) -> str:
    out = text
    for src, dst in SHIELD_TERMS.items():
        out = out.replace(src, dst)
    return out
\end{codeblock}

\begin{codeblock}[language=Python]{FastAPI endpoint sketch}
@router.get("/api/v1/alerts", response_model=AlertPage)
async def list_alerts(
    rule_name: str | None = None,
    severity: str | None = None,
    status: str | None = None,
    entity_type: str | None = None,
    entity_id: str | None = None,
    date_from: date | None = None,
    date_to: date | None = None,
    page: int = 1,
    page_size: int = 50,
):
    return await alert_query_service.fetch_page(...)
\end{codeblock}

\begin{codeblock}[language=Python]{DOCX generation (python-docx)}
from docx import Document

def build_complaint_docx(alert, explanation, output_path: str):
    doc = Document()
    doc.add_heading("Draft Complaint for eOtinish", level=1)
    doc.add_paragraph(f"Alert ID: {alert.alert_id}")
    doc.add_paragraph(f"Rule: {alert.rule_name} v{alert.rule_version}")
    doc.add_paragraph("Factual evidence:")
    for k, v in explanation["evidence"].items():
        doc.add_paragraph(f"- {k}: {v}")
    doc.add_paragraph(
      "Disclaimer: This report is an algorithmic risk assessment based on open data."
    )
    doc.save(output_path)
    return output_path
\end{codeblock}

\section{Phase 8: DevOps, Deployment, \& Human-in-the-Loop}

\subsection{Conceptual explanation (for beginners)}
\begin{conceptbox}
This solution is not a single application. It is a distributed data factory plus analytics and legal-output platform. Stable operations require service decomposition, infrastructure boundaries, and clear human checkpoints.

AI accelerates implementation, but several responsibilities remain human-critical:
\begin{itemize}[leftmargin=1.4em]
  \item anti-bot operational tactics,
  \item production reliability and incident response,
  \item risk-weight calibration governance,
  \item legal accountability and external communications.
\end{itemize}
\end{conceptbox}

\subsection{Technical design}
\textbf{Dockerized services:}
\begin{itemize}[leftmargin=1.4em]
  \item \texttt{api} (FastAPI presentation/query),
  \item \texttt{ingest-worker} (collectors),
  \item \texttt{normalize-worker} (validation + raw to core),
  \item \texttt{outbox-relay} (outbox to queue),
  \item \texttt{scoring-worker} (rule dispatcher),
  \item \texttt{doc-service} (PDF/DOCX generation),
  \item \texttt{postgres-master} (OLTP),
  \item \texttt{postgres-analytics} (OLAP + PostGIS/Timescale),
  \item \texttt{neo4j}, \texttt{redis}/\texttt{rabbitmq},
  \item \texttt{prometheus} + \texttt{grafana} + \texttt{loki}.
\end{itemize}

\textbf{Physical split topology:}
\begin{itemize}[leftmargin=1.4em]
  \item \textbf{Server A (Kazakhstan):} low-latency ingestion gateway, raw write path, minimal external exposure.
  \item \textbf{Server B (abroad):} heavy analytics, graph processing, document assembly.
  \item Private interconnect via WireGuard tunnel; no public database ports.
\end{itemize}

\textbf{Human-in-the-loop checkpoints:}
\begin{itemize}[leftmargin=1.4em]
  \item Data quality console for dead-letter triage,
  \item Manual whitelist/exception workflow for known mass-registration addresses,
  \item Rule tuning UI for weights/thresholds,
  \item Legal review queue before complaint submission.
\end{itemize}

\subsection{Key snippets}
\begin{codeblock}[language={}]{docker-compose (condensed)}
version: "3.9"
services:
  postgres-master:
    image: postgres:15
    environment:
      POSTGRES_DB: audit
      POSTGRES_USER: audit
      POSTGRES_PASSWORD: audit
    ports: ["5432:5432"]
    volumes: ["pg_master:/var/lib/postgresql/data"]

  postgres-analytics:
    image: timescale/timescaledb-ha:pg15
    environment:
      POSTGRES_DB: audit_olap
      POSTGRES_USER: audit
      POSTGRES_PASSWORD: audit
    ports: ["5433:5432"]
    volumes: ["pg_olap:/var/lib/postgresql/data"]

  neo4j:
    image: neo4j:5
    environment:
      NEO4J_AUTH: neo4j/password
    ports: ["7474:7474", "7687:7687"]

  redis:
    image: redis:7-alpine
    ports: ["6379:6379"]

  rabbitmq:
    image: rabbitmq:3-management
    ports: ["5672:5672", "15672:15672"]

  api:
    build: ./services/api
    depends_on: [postgres-analytics, redis]
    ports: ["8000:8000"]

  ingest-worker:
    build: ./services/ingest
    depends_on: [postgres-master, rabbitmq]

volumes:
  pg_master:
  pg_olap:
\end{codeblock}

\begin{codeblock}[language={}]{WireGuard concept}
# /etc/wireguard/wg0.conf (example structure)
[Interface]
Address = 10.8.0.1/24
PrivateKey = <serverA_private_key>
ListenPort = 51820

[Peer]
PublicKey = <serverB_public_key>
AllowedIPs = 10.8.0.2/32
Endpoint = <serverB_public_ip>:51820
PersistentKeepalive = 25
\end{codeblock}

\subsection{Final implementation guidance (practical sequence)}
\begin{enumerate}[leftmargin=1.5em]
  \item Stand up OLTP + OLAP + queue + Neo4j via Docker.
  \item Create schemas/tables, including \texttt{ops.outbox\_events} and \texttt{ops.dead\_letter}, first.
  \item Implement raw ingestion and watermarking.
  \item Add validation + idempotent upsert + outbox transaction.
  \item Implement outbox relay and change queue workers.
  \item Build scoring dispatcher and first four rules.
  \item Add legal-safe document generation (PDF + DOCX).
  \item Add observability, dead-letter console, and rule-tuning UI.
  \item Deploy split topology (KZ gateway + abroad analytics) over WireGuard.
  \item Run one-week pilot, calibrate, then backfill the 2024--2025 archive.
\end{enumerate}

\vspace{1.2em}
\begin{center}
\textcolor{MutedGray}{\rule{0.75\textwidth}{0.4pt}}
\end{center}
\begin{center}
\small\textcolor{MutedGray}{Algorithmic risk assessment based on open data; not a direct assertion of criminal guilt.}
\end{center}

\end{document}
